\section{Introduction}
\label{sec:intro}

A deployed LLM system makes many small decisions around each answer it writes.
A router decides which model should handle a prompt~\cite{chen2023frugalgpt,ong2024routellm}.
A semantic cache decides whether a stored answer can be served instead of a new one~\cite{bang2023gptcache,schroeder2025vcache}.
An agent decides which tool to call next, or whether to stop and reply~\cite{yao2024taubench}.
Each of these decisions has a small, fixed set of outcomes (which model, reuse or regenerate, which tool), yet the usual way to make one is to ask a generative model and parse its text: the system pays for a full autoregressive call when all it needs back is one choice among a handful of options.

A recent class of models targets exactly this gap.
These \emph{decision models}, which their vendors call ``System One'' models after the fast mode of dual-process theory~\cite{kahneman2011thinking}, are not autoregressive: given a state and a typed question (pick one of these options, score on this scale, yes or no), they return the answer with a probability for each option in a single pass.
The first commercial one, \jev~\cite{typesafe2026jev}, is priced at a small fraction of a frontier LLM per input token with no output cost, and is advertised as a drop-in replacement for an LLM wherever the output is a decision.
An open-weights reproduction, \kev~\cite{palmer2026kev}, followed within weeks.

Whether this replacement pays is an empirical systems question, and it has two parts that vendor numbers do not separate.
The first is quality: is the cheap decision as good as the expensive one?
The second is structure: what does the decision buy?
A cache guard's decision can remove a generation entirely, so a cheap, accurate guard saves nearly the whole call.
An agent's tool choice sits in front of a generation that happens anyway, because some model must still write the tool's arguments, so a separate decider adds a call instead of removing one.

We study both parts on one harness.
Every decider implements the same interface (a state, a set of options, a returned choice with a confidence), so the same code path scores a frontier LLM, a small hosted LLM, a self-hosted LLM, \jev, \kev, and trained classifiers.
We run them at three decision points on public datasets, with repeated calls, clustered confidence intervals, measured cost per call, and measured GPU energy for the self-hosted deciders.

Our findings are as follows.

\textbf{Cache guard: the decision model pays off.}
GSM-Plus~\cite{li2024gsmplus} rewrites grade-school math questions into variants: some only reword the question and keep its answer, others change a number and so change the answer.
Used as cache traffic, a new question and its nearest cached neighbour often differ by that one number, so embedding similarity cannot tell a safe reuse from a wrong one, and the best similarity cut-off, tuned on held-out queries, reuses almost nothing (66.1\% accuracy).
\jev\ reaches 93.3\%, within 2.6 points of the frontier model, at about 2\% of its price per decision.
A frontier model used as the guard is self-defeating: its decision costs more than the answer it might save.
On two datasets of real traffic, chat prompts and web search queries, \jev\ ties the similarity cut-off on chat, where similarity works, and beats it by 4.4 points on search, where it does not.
A cascade that uses the free similarity score for clear cases and calls \jev\ only when the score is ambiguous is best or tied for best on all three datasets.
(On the two traffic datasets every decider, the frontier model included, first stalled near 71\% accuracy; we traced this to errors in the datasets' labels, relabeled the pairs we use, and report results on the corrected labels, Section~\ref{sec:labels}.)

\textbf{Agent control: it does not.}
\taubench~\cite{barres2025tau2} simulates airline customer-service conversations in which an agent must look up and change bookings through tools while following a written policy.
At saved points in such conversations, \jev\ picks the same next action as a frontier model's successful run 81.4\% of the time, at $1/22$ of the frontier model's cost and latency.
Inside a live agent loop, however, the standard design, in which the generating model decides and acts in one call, is the cheapest of all and matches the best reward.
Splitting the decision out either costs more (the generator asked twice) or loses about five points of reward (\jev\ deciding), a gap at the edge of significance.

\textbf{Routing: a trained classifier wins.}
RouterBench~\cite{hu2024routerbench} records which prompts a cheap model (Mixtral-8x7B) and an expensive one (GPT-4) answered correctly.
Asked to send each prompt to one of the two, no general decider, frontier LLM included, routes usefully from the prompt alone.
A small classifier trained on the cheap model's record of right and wrong answers beats all of them, and even it recovers little.

\textbf{The pattern.}
A standalone decider is worth its call when two conditions hold: the decision avoids a generation, and the evidence it needs is visible in its input.
Cache guarding meets both. Agent control fails the first and routing fails the second.

We make two contributions:
(1) an independent measurement of the decision-model paradigm at three LLM decision points, in cost, latency and energy, against open baselines that run without any API; and
(2) a decider interface and harness that make a typed-output decision model interchangeable with chat LLMs and classifiers.
